\clearpage\pagestyle{otMain}\setlength{\otBlocEspacement}{2.5mm}\fontsize{11}{15}\selectfont\corrigesChapitre{5}{Expressions et distributivité}
\begin{blocLivre}{cor-5-1}\begin{corrige}{1}{Lire la structure}\demarche{Repérer l’opération extérieure.}
$5x+15$ est une somme de deux termes. $5(x+3)$ et $x(2x-1)$ sont des produits : leur dernière opération est une multiplication. Les parenthèses contiennent elles-mêmes une somme ou une différence.
\end{corrige}\end{blocLivre}
\begin{blocLivre}{cor-5-2}\begin{corrige}{2}{Substituer un négatif}\demarche{Remplacer chaque lettre avant de calculer.}
On remplace $x$ par $-3$ : $A=4\times(-3)-7=-12-7=\res{-19}$. Les parenthèses montrent le nombre remplacé.
\end{corrige}\end{blocLivre}
\begin{blocLivre}{cor-5-3}\begin{corrige}{3}{Réduire}\demarche{Regrouper les termes de même nature.}
Les termes en $x$ donnent $8x+2x=10x$ et les constantes $-3+7=4$. L’expression réduite est $\res{10x+4}$.
\end{corrige}\end{blocLivre}
\begin{blocLivre}{cor-5-4}\begin{corrige}{4}{Développer}\demarche{Multiplier tous les termes de la parenthèse.}
On distribue le facteur 4 aux deux termes : $4\times3x-4\times2=\res{12x-8}$.
\end{corrige}\end{blocLivre}
\begin{blocLivre}{cor-5-5}\begin{corrige}{5}{Factoriser}\demarche{Chercher le facteur commun aux deux termes.}
Le facteur commun est 6 : $18x+30=6\times3x+6\times5=\res{6(3x+5)}$. Redévelopper permet de vérifier.
\end{corrige}\end{blocLivre}
\begin{blocLivre}{cor-5-6}\begin{corrige}{6}{Distribuer un signe négatif}\demarche{Contrôler le signe de chaque produit.}
$-2(5x-4)=-10x+8$, puis $-10x+8+3x=\res{-7x+8}$. Le terme $+8$ vient de $(-2)\times(-4)$.
\end{corrige}\end{blocLivre}
\begin{blocLivre}{cor-5-7}\begin{corrige}{7}{Un facteur littéral}\demarche{Utiliser $x	imes x=x^2$.}
$3x\times2x=6x^2$ et $3x\times5=15x$. On obtient $6x^2+15x-4x=\res{6x^2+11x}$. Les termes en $x^2$ et en $x$ restent séparés.
\end{corrige}\end{blocLivre}
\begin{blocLivre}{cor-5-8}\begin{corrige}{8}{Factoriser avec une lettre}\demarche{Mettre en évidence le facteur numérique et littéral.}
Chaque terme contient $5x$ : $15x^2=5x\times3x$ et $-10x=5x\times(-2)$. Donc $\res{5x(3x-2)}$.
\end{corrige}\end{blocLivre}
\begin{blocLivre}{cor-5-9}\begin{corrige}{9}{Équivalence de programmes}\demarche{Traduire les deux programmes avec la même lettre.}
Pour un nombre $x$, A donne $3(x+4)=3x+12$ et B donne $3x+12$. Les deux expressions sont identiques pour tout $x$ : les programmes sont équivalents.
\end{corrige}\end{blocLivre}
\begin{blocLivre}{cor-5-10}\begin{corrige}{10}{Réparer une réduction}\demarche{Vérifier les termes avant de les regrouper.}
On ne regroupe pas des termes de natures différentes. La somme reste $4x+3x^2$. Pour $x=2$, elle vaut $8+12=\res{20}$, tandis que $7x^3$ vaudrait 56 : ce contre-exemple réfute l’égalité.
\end{corrige}\end{blocLivre}
\begin{blocLivre}{cor-5-11}\begin{corrige}{11}{Une aire composée}\demarche{Additionner les aires des deux rectangles.}
Par la formule du rectangle, $A=4(x+3)$. En partageant la largeur en $x$ et 3, on obtient deux aires $4x$ et 12. Ainsi $\res{A=4x+12\ \mathrm{cm}^2}$. Cette décomposition illustre la distributivité.
\label{fig-cor-5-11}\figureGeo[corrige]{aire-distrib}
\end{corrige}\end{blocLivre}
\begin{blocLivre}{cor-5-12}\begin{corrige}{12}{Un périmètre variable}\demarche{Compter deux fois chaque dimension.}
$P=2(2x+1)+2(x+3)=4x+2+2x+6=\res{6x+8}$ cm. Pour $x=4$, $P=24+8=\res{32}$ cm. Les côtés valent alors 9 et 7 cm, d’où $2(9+7)=32$.
\label{fig-cor-5-12}\figureGeo[corrige]{rectangle-litteral}
\end{corrige}\end{blocLivre}
\begin{blocLivre}{cor-5-13}\begin{corrige}{13}{Un calcul mental justifié}\demarche{Repérer un facteur commun avant de calculer.}
Le facteur commun est 27 : $39\times27+61\times27=(39+61)\times27=100\times27=\res{2700}$. On a utilisé la distributivité dans le sens de la factorisation.
\end{corrige}\end{blocLivre}
\begin{blocLivre}{cor-5-14}\begin{corrige}{14}{Deux entiers consécutifs}\demarche{Représenter deux entiers consécutifs avec une seule lettre.}
On appelle $n$ le premier entier ; le suivant est $n+1$. Leur somme est $n+(n+1)=2n+1$. Elle a donc la forme d’un entier impair, pour tout entier $n$. Un seul exemple n’aurait pas suffi à le démontrer.
\end{corrige}\end{blocLivre}
\begin{blocLivre}{cor-5-15}\begin{corrige}{15}{Une différence constante}\demarche{Soustraire des expressions entre parenthèses.}
A donne $5x-7$. B donne $5(x-2)+6=5x-4$. Leur différence est $(5x-7)-(5x-4)=\res{-3}$. A donne donc toujours 3 de moins que B, quel que soit $x$. La lettre disparaît dans la différence.
\end{corrige}\end{blocLivre}
